\documentclass[12pt]{article}
\usepackage[utf8]{inputenc}
\usepackage[T1]{fontenc}
\usepackage{hyperref}
\title{Woi Language Resources for K-12 Curriculum Development}
\author{Compiled by Claude Sonnet 5 (claude-sonnet-5) for Andr\'as Kornai}
\date{}
\begin{document}
\maketitle

\section{Introduction}
Woi (also spelled Wooi; ISO 639-3 wbw) is an Austronesian language of Yapen Island in Papua, Indonesia, spoken mainly in the villages of Wooi and Woinap and on Moisnum Island. Speaker estimates differ between sources, from about 1,600 to about 3,000. Searches found academic documentation but no school materials.

\section{Primary Educational Resources}
No primers, textbooks or school materials were found in the search results.

\section{Linguistic and Reference Materials}
\begin{itemize}
\item Sawaki, \emph{A grammar of Wooi: An Austronesian language of Yapen Island, Western New Guinea} (ANU PhD thesis). Its abstract describes the phonology and a basic SVO clause order: \url{https://openresearch-repository.anu.edu.au/items/3858e53a-82dd-498b-a7d0-e4684df9702b}
\item Full text of the thesis (PDF): \url{https://openresearch-repository.anu.edu.au/bitstream/1885/136851/1/Sawaki%20Thesis%202017.pdf}
\item Wikipedia overview of Woi: \url{https://en.wikipedia.org/wiki/Woi_language}
\item Omniglot page on Wooi: \url{https://www.omniglot.com/writing/wooi.htm}
\item Max Planck EVA numerals page for Wooi: \url{https://lingweb.eva.mpg.de/channumerals/Wooi.htm}
\end{itemize}

\section{Religious and Community Resources}
None found in the search results.

\section{Digital and Audio Resources}
\begin{itemize}
\item DoBeS documentation project page: \url{https://dobes.mpi.nl/?p=1391}
\item Documentation of Wooi, Center for Endangered Languages Documentation (CELD) at Universitas Papua, which grew out of the project: \url{https://celd.phil-fak.uni-koeln.de/celd/the-documentation-of-wooi}
\end{itemize}

\section{Curriculum Development Context}
The documentation archives and the Sawaki grammar are the only substantial sources found. Any curriculum would need new materials developed with the community and with CELD. No school programme was found.
\end{document}